\subsection{The precise limit assertion}

We fix a rational tangential base point of
\(U=\mathbb P^1_{\mathbb Q}\setminus\{0,1,\infty\}\).
Put \(G=G_{\mathbb Q}\), and let \(F_\ell\) be the free pro-\(\ell\)
geometric fundamental group on the meridians at zero and one.
All sections, cocycles, and group homomorphisms below are continuous.
Conjugation of sections always means geometric conjugation.

\begin{result}{A.1}{Proved}{Comparison of the three survival assertions}
For an ordinary completed pair \(a=(A,B)\), write \(S_\ell(a)\)
for existence of a section of the geometrically pro-\(\ell\) quotient
\(E_U^{(\ell)}\) whose abelian class is \(a_\ell\).
Write \(S_{r,\ell}(a)\) for its class-\(r\) version,
\(S_r(a)\) for existence in \(E_U/\Gamma_{r+1}\Pi\), and
\(S_{\rm nil}(a)\) for existence in \(E_U^{\rm nil}\).
Then
\[
 S_{\rm nil}(a)\ \Longleftrightarrow\
 \bigl(\forall\ell\;S_\ell(a)\bigr)
 \ \Longleftrightarrow\
 \bigl(\forall r\;S_r(a)\bigr).
\]
The pair is the same in every quantifier. Define
\[
\begin{split}
 (\mathrm{FS})_{\rm nil}:\quad&
 \forall a\in\mathcal G_{\rm ord},\quad
 S_{\rm nil}(a)\Longrightarrow |\operatorname{supp}(a)|<\infty,\\
 (\mathrm{FS})_\ell:\quad&
 \forall a\in\mathcal G_{\rm ord},\quad
 S_\ell(a)\Longrightarrow |\operatorname{supp}(a)|<\infty.
\end{split}
\]
For any one prime \(\ell\), \((\mathrm{FS})_\ell\) implies
\((\mathrm{FS})_{\rm nil}\). A counterexample to the latter would
be a counterexample to every \((\mathrm{FS})_\ell\), with the same pair.
Counterexamples chosen separately for the different primes do not furnish
that simultaneous counterexample.
\end{result}
\begin{proof}
Every finite nilpotent group is the direct product of its Sylow subgroups.
Consequently the universal pronilpotent quotient of \(\Pi\) is
\(\prod_\ell F_\ell\), and its quotient of class at most \(r\) is
\(\prod_\ell F_\ell/\Gamma_{r+1}F_\ell\).
The action defined by the fixed tangential section preserves these
characteristic factors. Thus the arithmetic quotients are the semidirect
products with \(G\), and their projection maps are continuous
surjective homomorphisms over \(G\). A section is a crossed
homomorphism to the corresponding geometric group.
A tuple of continuous crossed homomorphisms is a continuous crossed
homomorphism to the product, since continuity into a product is tested
coordinatewise. This proves the first equivalence and
\(S_r(a)\Longleftrightarrow\forall\ell\,S_{r,\ell}(a)\).

For fixed \(\ell\), fix an actual abelian cocycle representing
\(a_\ell\). Any section with this abelian cohomology class can be
geometrically conjugated to have this literal cocycle: lift the element
giving its abelian coboundary through the surjection of geometric groups.
The fixed-cocycle compactness theorem independently proved below shows
that the spaces at all integral classes are compact. If each is nonempty,
the equations for adjacent projection maps have the finite intersection
property: choose a section at the largest class occurring in a finite
list of equations and project it. Their common solution is a compatible
sequence. Its inverse limit is continuous, and
\(F_\ell=\varprojlim_r F_\ell/\Gamma_{r+1}F_\ell\), since every
finite \(\ell\)-group is nilpotent. This proves the second equivalence.
The remaining implications follow by applying the projections to the
same pair. No interchange of \(\forall\ell\) with \(\exists a\)
is involved.
\end{proof}



\begin{result}{A.2}{Open}{The main question and its single-prime versions}
We have neither proved \((\mathrm{FS})_{\rm nil}\) nor constructed
an ordinary infinite-support pair with a section of \(E_U^{\rm nil}\).
We have also not decided \((\mathrm{FS})_2\), \((\mathrm{FS})_3\),
or the assertion for arbitrary \(\ell\ge5\).
These are the exact remaining assertions in tasks A2 and A3 and Lead L1.
The fixed-pair equivalence in A.1 is a reduction to an integral class
obstruction for each ghost; it is not an obstruction itself.
At coefficient three, E.146--E.148 give an exact finite extension
criterion: the old-support plus-character Frobenius equations, with
every aligning frame retained. E.151 preserves an ordinary pair in
the full power completion along any actually chosen infinite path
through these systems. The missing assertion is existence or
impossibility of such a compatible path. E.149 shows that individual
frozen predecessors can fail, so an induction asserting that every
node extends is false. This conditional construction concerns one
specified support pattern; failure of that entire pattern would not
by itself prove the general finite-support assertion.
At coefficient two, E.155--E.159 remove the real obstruction for
an actually boundary-matched ordinary predecessor and give an exact
finite-parameter test for preserving every off-diagonal support match
in the stated support pattern. Its passing infinite branch is not
constructed. E.171--E.174 do construct one ordinary infinite-support
pair with a single integral class-two section matching every support
point away from its coefficient prime, simultaneously for all
coefficient primes; that section has an unrestricted global class-three
lift. These assertions establish nonemptiness of the necessary
class-two matching locus. They neither preserve point matching at
class three nor supply any all-depth section, and they concern a
different family from the original A.44 recursion.
\end{result}

\begin{result}{A.3}{Proved}{Why finite-ramification pro-prime constructions cannot be ghosts}
For an ordinary pair, its valuations are actual integers. If its
\(\ell\)-primary abelian cocycle is unramified outside a finite set,
then the pair has finite prime support and is a rational unit-equation
pair. In particular, equality of its \(\ell\)-primary pair with that
of a rational point, even at one coefficient prime, implies finite support.
For a ghost, the ordinary valuation vectors outside every finite set
span a rank-two subgroup of \(\mathbb Z^2\), and generate an open
\(\mathbb Z_\ell\)-submodule of \(\mathbb Z_\ell^2\).
Its pro-\(\ell\) abelian pair is not cuspidal.
This confirms these assertions of GT\_TOWER\_CHECK C.1.2.
\end{result}
\begin{proof}
At a finite place \(q\ne\ell\), the restriction to tame
\(\ell\)-inertia of the Kummer class of \(u\in\mathbb Q_q^\times\)
is \(v_q(u)\) times the tame character into \(\mathbb Z_\ell(1)\).
This follows by writing \(u=q^m v\), adjoining compatible
\(\ell\)-power roots of the uniformizer, and observing that units have
unramified \(\ell\)-power Kummer classes. A nonzero integer is nonzero
in \(\mathbb Z_\ell\). Thus every nonzero valuation of either ordinary
coordinate at \(q\ne\ell\) ramifies that abelian class. The remaining
place \(\ell\) adds at most one support prime. The finite-support
reconstruction is independently proved in E.101. Equality with a rational
\(\ell\)-primary pair makes all these integer valuations equal to the
rational valuations, hence confines the support to a finite set.

The three possible nonzero ordinary valuation patterns lie on the
distinct rational lines \((m,0)\), \((0,m)\), and \((m,m)\).
If their span outside a finite set had rank at most one, only one of
these lines could occur there. Accordingly one of \(B,A,A/B\)
would have finite support, and E.101 would force the pair to be
rational, a contradiction. Two independent integer vectors have
nonzero determinant and generate an open submodule over every
\(\mathbb Z_\ell\). Finally a zero-cusp, one-cusp, or infinity-cusp
abelian pair has respectively \(B=1\), \(A=1\), or \(B/A=-1\)
in the pro-\(\ell\) completion. The corresponding ordinary
integer valuations vanish, again giving finite support for one
coordinate or the ratio and contradicting ghosthood.
\end{proof}

\subsection{The consequence for full sections is stronger than stated in the prompt}

\begin{result}{A.4}{Known}{Rational abelian pair and ordinary local full sections determine the full section}
Let \(s:G\to E_U\) be a section of the full profinite arithmetic
fundamental group, whose restriction at every finite place is the section
of an ordinary point \(x_q\in U(\mathbb Q_q)\). If its completed
abelian pair is \((\kappa(t),\kappa(1-t))\) for
\(t\in\mathbb Q\setminus\{0,1\}\), then \([s]=[s_t]\).
The derivation uses \citet[Corollary 5, p.~5]{Stix2015}, or equivalently
his Corollary 6, p.~6, after checking adelicity; it does not use the
attached unrefereed local theorem.
\end{result}
\begin{proof}
For every finite \(q\), the map
\(\mathbb Q_q^\times\to\varprojlim_n
\mathbb Q_q^\times/(\mathbb Q_q^\times)^n\) is injective.
Indeed its valuation factor is \(\mathbb Z\), its roots-of-unity
factor is finite, and its principal-unit factor has an open subgroup
isomorphic to the additive group \(\mathbb Z_q\). The intersection
of the subgroups of all powers in each factor is trivial. Consequently
\(\kappa(x_q)=\kappa(t)\) implies \(x_q=t\).
These local points are integral on \(U\) outside the finite set of
primes dividing the numerator or denominator of \(t(1-t)\).

At the real place the real section theorem identifies full section
classes with connected components. Its three ordinary components have
Kummer sign vectors \((- ,+),(+,+),(+,-)\). The rational abelian pair
therefore gives the component of \(t\); the exact real result and its
source are recorded in Part P. Thus \(s\) is an adelic section in
Stix's sense. The regular unit \(t\) on the tripod is nonconstant,
so all hypotheses of his Corollary 6 hold. It follows that \(s\) is
the section of a rational point. Its first Kummer coordinate identifies
that point as \(t\). Alternatively, its local points have the finite
geometric support \(\{t\}\), so Corollary 5 gives the same result.
This is finite geometric support of local points, a different notion
from finite prime support of their valuations.
\end{proof}

\begin{result}{A.5}{Proved}{The exact conditional implication of a positive or negative limit answer}
If \((\mathrm{FS})_{\rm nil}\) holds, every section of \(E_U\to G\)
whose local sections at finite primes are ordinary is a rational-point
section. Thus a positive answer would prove the locally ordinary case of
Conjecture A, not merely reduce it to an unresolved rational-abelian-pair
case. The twenty-four unrestricted class-two lifts are not a remaining
ambiguity under these hypotheses. A negative answer to
\((\mathrm{FS})_{\rm nil}\) alone would give a non-point section of
\(E_U^{\rm nil}\); to contradict the full section conjecture one would
still have to lift it to \(E_U\).
\end{result}
\begin{proof}
The quotient homomorphism \(E_U\twoheadrightarrow E_U^{\rm nil}\)
is continuous and lies over the identity of \(G\). Composing it with
\(s\) preserves its abelian pair. Local point sections also project
functorially, so the pair is ordinary. The hypothetical finite-support
assertion and E.101 give a rational pair. Apply A.4 to the original
full section. This proves the affirmative implication without asserting
that a pronilpotent section lifts back to a full section.

Conversely, an infinite-support pair cannot be the pair of a rational
point by E.101. Its pronilpotent section, if constructed, would therefore
be non-point. The quotient map on groups supplies a map on section
classes in this direction only; it supplies no section of the map on
section sets. The separate lift to the full group is an additional
embedding problem. Finally, A.29 below shows explicitly that the local
class-two comparisons at \((q,\ell)=(3,2),(2,3)\) already distinguish
all twenty-four class-two twists.
\end{proof}

\begin{result}{A.6}{Proved}{The obstruction at an arbitrary integral class is a class of a specified lift}
Fix \(\ell\), a literal class-\(r\) cocycle \(c\), and the actual
arithmetic action. A continuous Hall-coordinate section of
\(N_{r+1,\ell}\twoheadrightarrow N_{r,\ell}\) gives a lift
\(\widetilde c\) as a map of spaces. The formula
\[
 o_c(g,h)=\widetilde c(g)\,g\widetilde c(h)\,
                 \widetilde c(gh)^{-1}
\]
defines a continuous two-cocycle in
\(L_{r+1,\ell}=\mathbb Z_\ell(r+1)^{b_{r+1}}\).
Its class vanishes precisely when \(c\) lifts. It depends on the
class-\(r\) section as well as on the abelian pair. Its projections to
finite coefficients need not determine whether this continuous class
vanishes.
\end{result}
\begin{proof}
The quotient of the displayed expression is one by the cocycle identity
for \(c\), hence it lies in the central kernel. Associating a product
of three lifted factors in the two possible ways is exactly the
two-cocycle identity for \(o_c\). Inner conjugation from
\(N_{r,\ell}\) acts trivially on the central layer, so its action is
the original cyclotomic weight \(r+1\) action. Replacing
\(\widetilde c\) by a central continuous cochain changes \(o_c\)
by its coboundary. A primitive gives the required corrected crossed
homomorphism; a lift supplies that primitive by taking its difference
from \(\widetilde c\). The continuous-versus-finite coefficient
qualification is the exact derived-limit sequence proved and applied
below. This constructs the obstruction map, but does not evaluate it
for every ghost at every class.
\end{proof}

\begin{result}{A.7}{Open}{The unsolved evaluation problem in A4 and L4}
No theorem here identifies a first nonzero higher obstruction for every
ordinary ghost. In particular, the existence of nonzero locally trivial
continuous cohomology classes does not prove that an obstruction of a
given ghost is one of them. The exact class-three residue test below
separates local vanishing from the additional convergence of primitives.
This supplies a precise gate, but neither proves the nonvanishing
claimed in L4 nor constructs an all-depth survivor.
\end{result}

\begin{result}{A.8}{Proved}{Compatibility, prescribed symmetries, and the ambiguity of a lift}
Let \(X\) be a compact space of candidates, and let
\(X_i\subseteq X\) be closed subsets expressing specified finite
conditions. If every finite collection of these conditions has a
simultaneous solution in \(X\), they have a common solution.
Consequently any purported example in which all finite cofinal tests
of one fixed closed lifting-and-symmetry problem succeed, but the
problem has no solution, must fail one of these hypotheses.
For sections, a chosen lift through a central kernel \(T\), if it
exists, leaves the torsor under \(H^1(G,T)\) undetermined.
This is the relevant test of the comparisons proposed in R9.
\end{result}
\begin{proof}
The first assertion is the finite-intersection characterization of
compactness. If the condition is that a lifted automorphism be an
involution, its square is a continuous function into the relevant
profinite group. Vanishing of the square is equivalent to vanishing
in every member of a cofinal system of finite quotients, with a
single compatible automorphism as variable. A tower which is not
cofinal, changing prescribed frames, or choices whose projection
values change the original problem do not give these hypotheses.

For the section assertion, choose one lift \(s\) through a central
kernel. Any other lift with identical image is uniquely
\(s'(g)=b(g)s(g)\), and the homomorphism identity is precisely the
one-cocycle identity for \(b:G\to T\). Conjugation by \(T\)
changes \(b\) by a coboundary. Thus the quotient of the literal
lift fibre by these conjugations is a torsor under \(H^1(G,T)\).
Further geometric conjugations may act through stabilizers of the
prescribed quotient section; they must be specified separately.
The compactness proof A.20 fixes the literal abelian cocycle and
retains all the higher cocycles as its parameters, so it makes no
claim of a distinguished lift. Nothing in this argument uses the
arithmetic assertions of LM4\_SUCCESSOR A.79, A.81, or B.42--B.45
as established facts; their remaining audit status is recorded in
Part E. A conditional analogy with those assertions cannot replace
the hypotheses of this proof.
\end{proof}

\begin{result}{A.9}{Proved}{All-depth survival alone does not bound ramification}
Choose distinct rational primes \(q_j\), and put
\(A=\prod_{j\ge0}q_j^{(j+4)!}\) in the power completion.
There is an actual full profinite cuspidal section of the tripod with
abelian pair \((A,1)\). It has infinite support and ramifies at infinitely
many primes in every coefficient component. Its pair is not ordinary.
Thus finite ramification cannot follow from all-depth lifting alone;
any proof of A2 must use the ordinary local unit equation.
\end{result}
\begin{proof}
Fix compatible roots of each \(q_j\) and use the Kummer cocycle
\(a=\sum_j(j+4)!\kappa(q_j)\) with values in
\(\widehat{\mathbb Z}(1)\). Its tails vanish uniformly at every fixed
finite coefficient level, so it is a continuous cocycle. At the tangent
\(\overrightarrow{01}\), the actual action on the zero meridian is
\(g(x)=x^{\chi(g)}\). Therefore
\(s(g)=x^{a(g)}s_0(g)\) is a continuous homomorphism splitting
\(E_U\to G_{\mathbb Q}\); it lies in the zero-cusp decomposition
group. Its abelianization is the displayed pair. Every chosen prime has
nonzero ordinary valuation \((j+4)!\), and the inertia calculation of
A.3 proves the ramification assertion at each coefficient prime.
If this pair were ordinary at a finite prime, local full Kummer
injectivity would force \(1-x_q=1\), hence \(x_q=0\), which is not
in \(U\). This establishes the stated scope of the example.
\end{proof}
